\section{Block extensions and stable tameness}\label{sec:blocks}

We now consider higher-dimensional extensions, separating the base variables from the last three variables. Let $n\geq6$, put $m=n-6$, and write
\[
 \iota=(x_1,x_2,x_3),\qquad \varphi=(x_4,\ldots,x_{n-3}),\qquad
 \psi=(x_{n-2},x_{n-1},x_n).
\]
The tuple $\varphi$ is empty when $n=6$. Set $\zeta=(\iota,\varphi)$ and $R=K[\zeta]$, $L=\Frac R$. We consider maps of the form
\begin{equation}\label{eq:block-map}
 \widetilde H(\iota,\varphi,\psi)=\bigl(A(\iota),C(\iota),B(\zeta,\psi)\bigr),
\end{equation}
where $A\in K[\iota]^3$, $C\in K[\iota]^m$, and $B\in R[\psi]^3$.

For a field $K$, we use the notation
\begin{equation}\label{eq:D-family}
 \begin{aligned}
 \Dcal_{\vartheta_1,\vartheta_2,h,d_1,d_2}(r_1,r_2,r_3)
 =\bigl(&\vartheta_2h(\vartheta_1r_1+\vartheta_2r_2)+d_1,\\
       &-\vartheta_1h(\vartheta_1r_1+\vartheta_2r_2)+d_2,\ 0\bigr),
 \end{aligned}
\end{equation}
where $\vartheta_1,\vartheta_2,d_1,d_2\in K[r_3]$, $(\vartheta_1,\vartheta_2)=K[r_3]$, and $h\in K[r_3][t]$ with $h(0)=0$. We retain $\Ncal_P$ for the independent family~\eqref{eq:canonical}.

\subsection{Normal forms of the two blocks}
\begin{theorem}[Blockwise normal forms]\label{thm2}
Suppose that $\widetilde H$ has the form~\eqref{eq:block-map}, $\widetilde H(0)=0$, and $J\widetilde H$ is nilpotent. Then the following statements hold.
\begin{enumerate}
\item There is $S\in\GL_3(K)$ such that $S\circ A\circ S^{-1}$ is either a map of the form~\eqref{eq:D-family} over $K$, or $\Ncal_P$ for a nonconstant $P\in K[t]$ with $P(0)=0$.
\item Put $\chi(\zeta)=B(\zeta,0)$ and $B^0(\zeta,\psi)=B(\zeta,\psi)-\chi(\zeta)$. Viewed as a polynomial map in $\psi$ over $L$, the normalized map $B^0$ admits $T\in\GL_3(L)$ such that
\begin{equation}\label{eq:fiber-normal-form}
 T B^0(\zeta,T^{-1}\psi)=\Dcal_{\upsilon_1,\upsilon_2,f,\omega_1,\omega_2}(\psi)
 \quad\text{or}\quad \Ncal_Q(\psi).
\end{equation}
In the first case, $\upsilon_1,\upsilon_2,\omega_1,\omega_2\in L[\psi_3]$ and $f\in L[\psi_3][t]$, with $(\upsilon_1,\upsilon_2)=L[\psi_3]$ and $f(0)=0$. In the second case, $Q\in L[t]\setminus L$ and $Q(0)=0$.
\end{enumerate}
The constant conjugation $\widehat S=\diag(S,I_m,I_3)$ of the full map is
\begin{equation}\label{eq:whole-base-conjugation}
 \widehat S\circ\widetilde H\circ\widehat S^{-1}
 =\bigl(SA(S^{-1}\iota),\ C(S^{-1}\iota),\ B(S^{-1}\iota,\varphi,\psi)\bigr).
\end{equation}
The transformation in~\eqref{eq:fiber-normal-form} changes the three polynomial variables over $L$, while keeping the base parameters fixed.
\end{theorem}

\begin{proof}
The Jacobian is block lower triangular:
\begin{equation}\label{eq:block-jacobian}
 J\widetilde H=
 \begin{pmatrix}
 J_\iota A&0&0\\
 J_\iota C&0&0\\
 J_\iota B&J_\varphi B&J_\psi B
 \end{pmatrix}.
\end{equation}
Its diagonal blocks $J_\iota A$ and $J_\psi B$ are nilpotent. Since $A(0)=0$, Theorem~\ref{thm:classification} applies if the components of $A$ are independent. Otherwise, a constant conjugation makes its third component vanish, and Lemma~\ref{lem:planar-form}, with $\iota_3$ as the coefficient variable, gives~\eqref{eq:D-family}. This proves (i), and direct substitution gives~\eqref{eq:whole-base-conjugation}.

Over $L$, the map $B^0$ satisfies $B^0(0)=0$ and $J_\psi B^0=J_\psi B$. Again distinguish the cases of dependent and independent components, now over the characteristic-zero field $L$. The independent case follows from Theorem~\ref{thm:classification}. In the dependent case, a matrix in $\GL_3(L)$ makes one component zero, and Lemma~\ref{lem:planar-form} applies over $L[\psi_3]$. This proves (ii).
\end{proof}

\begin{remark}\label{rem:block-interpretation}
The theorem concerns the two three-variable diagonal blocks. A matrix with entries in $K(\zeta)$ is constant for differentiation in $\psi$, but need not be constant for the full Jacobian. Thus~\eqref{eq:fiber-normal-form} cannot be substituted into~\eqref{eq:constant-conjugacy} for the full map. In particular, the notation $T\circ\widetilde H\circ T^{-1}$ with $T\in\GL_n(K(\zeta))$ does not by itself define a constant linear conjugation of $\widetilde H$. Formula~\eqref{eq:whole-base-conjugation} also shows why the middle block changes when the first three source coordinates change.
\end{remark}

The subtraction of $\chi(\zeta)$ is necessary even when $\widetilde H(0)=0$.

\begin{example}\label{ex:fiber-constant}
Take $n=6$, $A=0$, and
\[
 B(\zeta,\psi)=\bigl(\varepsilon,\ \psi_3-2\psi_1\varepsilon,\ \varepsilon^2+\zeta_1\bigr),\qquad \varepsilon=\psi_2+\psi_1^2.
\]
The full map $(0,0,0,B)$ vanishes at the origin and has nilpotent Jacobian by~\eqref{eq:block-jacobian} and~\eqref{eq:canonical-J}. Over $L=K(\zeta_1,\zeta_2,\zeta_3)$, its last three components are linearly independent. Their relation ideal is
\[
 (W-U^2-\zeta_1)\subset L[U,V,W],
\]
because $\varepsilon$ and $\psi_3-2\psi_1\varepsilon$ are algebraically independent. In particular, every nonzero relation of total degree at most two has a nonzero constant term.

A dependent form~\eqref{eq:D-family} has a nonzero linear relation. A canonical form $\Ncal_Q$ has the quadratic relation $W-U^2=0$, whose constant term is zero, even if one allows $Q(0)\neq0$. Constant linear target changes preserve the constant term of a relation, and source changes preserve the relation ideal. Therefore $B$ is not linearly conjugate over $L$ to either displayed family. After subtracting $B(\zeta,0)=(0,0,\zeta_1)$, the resulting map is exactly $\Ncal_t$.
\end{example}

\subsection{Automorphisms over the base ring}
The normal form over the fraction field alone does not establish stable tameness over $R$. We first prove polynomial invertibility over $R$.

\begin{lemma}\label{lem:inverse-descent}
Let $R$ be a domain, $L=\Frac R$, and let $G\in R[\psi_1,\ldots,\psi_r]^r$ satisfy $G(0)=0$ and $JG(0)\in\GL_r(R)$. If $G\in\GA_r(L)$, then $G\in\GA_r(R)$.
\end{lemma}

\begin{proof}
The formal inverse theorem over a ring gives a unique inverse $Q\in R[[\psi]]^r$ with $Q(0)=0$. It can be constructed degree by degree: after the terms below degree $d$ are known, the degree-$d$ correction is obtained by multiplying the degree-$d$ error by $JG(0)^{-1}$. Thus all coefficients lie in $R$.

The polynomial inverse of $G$ over $L$ fixes the origin. As an element of $L[[\psi]]^r$, it equals $Q$ by uniqueness. Hence $Q$ has only finitely many nonzero coefficients and belongs to $R[\psi]^r$. Both inverse identities hold over $R$, either by the formal construction or by the injection $R[\psi]\hookrightarrow L[\psi]$.
\end{proof}

\begin{proposition}\label{prop:ring-nilpotent}
Let $R=K[\zeta_1,\ldots,\zeta_r]$, and let $B\in R[\psi_1,\psi_2,\psi_3]^3$ have nilpotent Jacobian $J_\psi B$. Then
\[
 \Theta(\psi)=\psi+B(\zeta,\psi)
\]
is a polynomial automorphism over $R$ and is stably tame over $R$.
\end{proposition}

\begin{proof}
Put $\chi=B(\zeta,0)$ and $G(\psi)=\psi+B(\zeta,\psi)-\chi$. Then $G(0)=0$. The matrix $N=J_\psi B(\zeta,0)$ satisfies $N^3=0$, so
\[
 JG(0)=I+N\in\GL_3(R),\qquad (I+N)^{-1}=I-N+N^2.
\]
Over $L=\Frac R$, Theorem~\ref{thm1} implies that $G$ is a polynomial automorphism. Lemma~\ref{lem:inverse-descent} gives $G\in\GA_3(R)$. Since $\Theta=\trans{\chi}\circ G$, we also have $\Theta\in\GA_3(R)$.

For every prime $\mathfrak p\in\Spec R$, let $\kappa(\mathfrak p)$ be the residue field. It contains $K$, so it has characteristic zero. Nilpotence is expressed by a polynomial matrix identity and is therefore preserved under specialization. Thus
\[
 \Theta_{\mathfrak p}=\psi+B_{\mathfrak p}(\psi)
 \in\TA_3(\kappa(\mathfrak p))
\]
by Theorem~\ref{thm1}. The ring $R$ is regular. The residue-field criterion~\cite[Theorem~4.12]{BEW} therefore implies that $\Theta$ is stably tame over $R$.
\end{proof}

\subsection{Stable tameness of the full map}
\begin{theorem}[Stable tameness of block extensions]\label{thm3}
Let $n\geq6$ and let $\widetilde H\in K[x_1,\ldots,x_n]^n$ have nilpotent Jacobian $J\widetilde H$ and satisfy
\[
 H_i\in K[x_1,x_2,x_3]\qquad(1\leq i\leq n-3).
\]
Then $F=\id+\widetilde H$ is a polynomial automorphism and is stably tame over $K$. In particular, this holds for the maps of Theorem~\ref{thm2}.
\end{theorem}

\begin{proof}
Write $F$ in block form:
\begin{equation}\label{eq:block-F}
 \begin{aligned}
 F(\zeta,\psi)&=\bigl(F_0(\zeta),\Theta_{\zeta}(\psi)\bigr),\\
 F_0(\iota,\varphi)&=(\iota+A(\iota),\varphi+C(\iota)),\qquad \Theta_{\zeta}(\psi)=\psi+B(\zeta,\psi).
 \end{aligned}
\end{equation}
The diagonal blocks in~\eqref{eq:block-jacobian} are nilpotent. By Theorem~\ref{thm1}, $g(\iota)=\iota+A(\iota)$ is tame over $K$. The base automorphism factors as
\[
 F_0=(\iota,\varphi+C(g^{-1}(\iota)))\circ(g(\iota),\varphi).
\]
The first factor is a product of elementary maps, one for each middle coordinate. Hence $F_0$ is tame. For $m=0$, this factor is the identity.

By Proposition~\ref{prop:ring-nilpotent}, $\Theta\in\GA_3(R)$ is stably tame over $R=K[\zeta]$. Thus $(\Theta,\psi_4,\ldots,\psi_{3+s})\in\TA_{3+s}(R)$ for some $s\geq0$. An elementary factor over $R$ is elementary over $K$ after the base variables $\zeta$ are included and held fixed. To treat the linear factors, use Suslin's theorem in the form
\begin{equation}\label{eq:Suslin}
 \GL_l(K[\zeta])=\langle E_l(K[\zeta]),\GL_l(K)\rangle,
 \qquad l\geq3;
\end{equation}
see~\cite[Theorem~3.23]{BEW}. Each elementary matrix in~\eqref{eq:Suslin} lifts to an elementary polynomial automorphism over $K$, while a matrix in $\GL_l(K)$ is linear over $K$. It follows that $\Psi(\zeta,\psi)=(\zeta,\Theta_{\zeta}(\psi))$ is stably tame over $K$.

Finally, $F=(F_0(\zeta),\psi)\circ\Psi$. The first factor is tame and the second is stably tame. Their composition is tame after a common stabilization. This also establishes polynomial invertibility. Explicitly, for a target $(\varpi,z)$ one first sets $\zeta=F_0^{-1}(\varpi)$ and then $\psi=\Theta_{\zeta}^{-1}(z)$; the coefficients of $\Theta^{-1}$ lie in $R$ by Proposition~\ref{prop:ring-nilpotent}.
\end{proof}

\begin{remark}\label{rem:block-scope}
Normalization of $\widetilde H$ is not needed in Theorem~\ref{thm3}, because Theorem~\ref{thm1} allows arbitrary constant terms and Proposition~\ref{prop:ring-nilpotent} explicitly subtracts the relative constant term. The argument gives no uniform bound on the number of stabilizing variables. It also does not assert tameness of the full map without stabilization.
\end{remark}
